\section{Risk Management and Execution}

\subsection{Risk Management Is the Constraint, Not the Decoration}

The risk-management book is clear about a principle that should remain structurally visible in the repo:
\begin{quote}
Trading should be adapted to risk management, not risk management adapted to the trader's impulse of the day.
\end{quote}

That principle matters even more in microstructure trading because:
\begin{itemize}
  \item decisions happen fast;
  \item many opportunities appear in clusters;
  \item execution friction can hide real risk;
  \item and the psychological temptation to overtrade is unusually strong.
\end{itemize}

\subsection{Per-Trade Risk and Per-Day Risk}

The materials suggest a layered loss budget:
\begin{itemize}
  \item define a daily acceptable loss based on real emotional and financial tolerance;
  \item derive per-trade risk from expected trading frequency;
  \item refuse trades whose natural stop does not fit that budget;
  \item stop the session once the daily limit is hit.
\end{itemize}

This is not presented as optional discipline.
It is the base operating system of the trader.

\subsection{Why Daily Stop Rules Matter So Much Here}

Intraday microstructure trading generates many local opportunities and many opportunities to rationalize.
The source material explicitly warns about revenge trading, loosening rules after a loss, and increasing size just because small planned profits feel emotionally unsatisfying.

That means the repository's eventual paper/live testing environment should include hard session controls such as:
\begin{itemize}
  \item daily realized loss limit;
  \item daily number-of-trades guardrail;
  \item cool-down after a string of failed setups;
  \item and possibly setup-specific disable flags once performance degrades intraday.
\end{itemize}

\subsection{Stop Logic}

In the discretionary material, a stop is not a random distance.
It is a statement about why the trade is wrong.

Examples:
\begin{itemize}
  \item bounce trade: the defending participant is gone or overwhelmed;
  \item breakout trade: the market failed to accept beyond the trigger;
  \item iceberg trade: hidden defense did not exist or was exhausted;
  \item robot trade: the robot stopped working or got absorbed;
  \item spring reversal: the opposing anchor did not produce the expected response.
\end{itemize}

Research implication: stop placement should remain tied to event logic, not only to volatility normalization.

\subsection{Take-Profit, Partials, and Adds}

The materials favor practical exit management over heroic forecasting.
Several themes repeat:
\begin{itemize}
  \item take partial profit at the first reliable objective;
  \item let the remainder continue only if the move still behaves well;
  \item use adds selectively and only within the original risk framework;
  \item do not add just because the initial entry feels emotionally uncomfortable.
\end{itemize}

That suggests future execution research should separate:
\begin{itemize}
  \item entry quality,
  \item management quality,
  \item and final exit quality.
\end{itemize}

\subsection{Execution Tradeoffs: Market Versus Limit}

\begin{longtable}{p{0.18\linewidth}p{0.36\linewidth}p{0.38\linewidth}}
\toprule
\textbf{Order type} & \textbf{Advantages} & \textbf{Risks and tradeoffs} \\
\midrule
Market order & Priority of execution, simplicity during fast moves, useful when timing is more important than a tiny price improvement & Slippage, unfavorable fills during volatility, and potential to pay the spread right when liquidity is poor. \\
Limit order & Better price control, potential maker economics, useful when leaning on a level or joining a defended participant & Non-fill risk, queue risk, adverse selection, and false confidence from a price that looks attractive but never actually fills. \\
\bottomrule
\end{longtable}

In this trading style, the execution decision is part of the strategy, not merely an implementation detail.

\subsection{Queue Position and Fill Uncertainty}

The source material does not formalize queue math, but the logic is implicit.
If many traders are leaning on the same visible density, queue priority matters.
If the trader expects to use limit orders near defended levels, the eventual live system should track at least:
\begin{itemize}
  \item whether the order was passive or aggressive;
  \item whether it received a full, partial, or no fill;
  \item how far price traveled without filling;
  \item and whether execution style degraded or improved realized expectancy.
\end{itemize}

\subsection{Low-Liquidity Traps}

Thin books create a nasty illusion:
\begin{itemize}
  \item stops look small on paper;
  \item visible densities look large;
  \item but fills become uncertain and slippage becomes nonlinear.
\end{itemize}

This is one reason the source materials emphasize that ``large'' is instrument-specific and day-specific.
Research code should therefore treat relative size and local liquidity context as more important than raw visible volume alone.

\subsection{Why Discretionary Traders Overtrade}

The materials imply several reasons:
\begin{itemize}
  \item too many visually tempting local patterns;
  \item boredom during slower periods;
  \item frustration after missed moves;
  \item desire to recover losses quickly;
  \item emotional discomfort with small planned profits;
  \item and confusion between action and edge.
\end{itemize}

Automation can help here, not by removing judgment entirely, but by enforcing:
\begin{itemize}
  \item session-level limits;
  \item setup definitions;
  \item entry gating;
  \item and structured review.
\end{itemize}

\subsection{Crypto Versus Russian Market Execution Differences}

\subsubsection{Crypto}

Crypto offers:
\begin{itemize}
  \item longer trading hours;
  \item high public visibility of order books and trades;
  \item fast experimentation and easier continuous logging;
  \item but also exchange-specific quirks, variable liquidity quality, and event-driven jumps.
\end{itemize}

\subsubsection{Russian Market / T-Invest / MOEX}

Russian-market execution carries different constraints:
\begin{itemize}
  \item session structure matters more explicitly;
  \item broker and infrastructure path matter more;
  \item some instruments may be heavily influenced by local participant behavior;
  \item and the repository's historical archive is particularly limited because it mainly reflects order-book data without complete trade flow.
\end{itemize}

\subsection{Operational Risk in Prop-Style Environments}

The MOEX prop book adds another layer:
\begin{itemize}
  \item commission burden matters materially;
  \item infrastructure quality matters;
  \item colocation matters for serious intraday operations;
  \item integrated risk management is a business necessity, not only a trader preference.
\end{itemize}

For the repository direction, this means architecture choices should already anticipate:
\begin{itemize}
  \item execution telemetry;
  \item per-session and per-strategy controls;
  \item and later integration between signal generation and hard risk guards.
\end{itemize}
